% GENERATED FILE. Edit canonical lessons, then run npm run generate:books.
% canonical-source-hash: fnv1a64:7528e9a44f24219a
% canonical-lessons: MW-C244-jhagro, MW-C244-larno, MW-C244-jhuth, MW-C244-dos, MW-C244-dhamki

\chapter{A Quarrel Next Door}
\label{ch:mw-a-quarrel-next-door}
\hlchaptermodality{244}

\begin{chapteropening}
\textbf{By the end of this chapter:} \emph{I can say who quarrelled and whose fault it was.}

Complete the last lesson of chapter 244: I can say who quarrelled and whose fault it was.
\end{chapteropening}

\section[\texorpdfstring{\mw{झगड़ो} (jhagṛo)}{झगड़ो (jhagṛo)}]{\mw{झगड़ो} (jhagṛo) — a quarrel}

\label{lesson:MW-C244-jhagro}



\begin{quote}
Before the new one: say the Marwadi for to go missing, then the Marwadi for to cheat.
\end{quote}

\subsection*{You'll want to know: \mw{झगड़ो}}
\textbf{\mw{झगड़ो}} — \emph{jhagṛo} — \textquotedblleft{}a quarrel\textquotedblright{}.

\begin{quote}
\mw{झगड़ो} \mw{मत} \mw{करो।}
\end{quote}

\begin{quote}
\emph{jhagṛo mat karo.}
\end{quote}

\begin{quote}
Don't quarrel.
\end{quote}

A masculine noun.

\begin{cousinweb}
\textbf{\mw{झगड़ो}} (\emph{jhagṛo}) is inherited from a reconstructed \emph{*jhagaḍ}, \textquotedblleft{}quarrel\textquotedblright{}; CDIAL lists Marwari \emph{jhagṛo} under it (CDIAL 5321).
\end{cousinweb}

\begin{grammarlens}[title={What you've built}]
One more word for saying what happened.
\end{grammarlens}

\subsection*{Your turn}
\begin{itemize}
\raggedright
  \item \emph{Say it:} \emph{jhagṛo}
  \item \emph{Say it:} \emph{jhagṛo}, once more
  \item \emph{Say it:} \emph{jhagṛo mat karo.} — Don't quarrel.
  \item \emph{Recall:} say the Marwadi for to go missing, then the Marwadi for to cheat, then say \emph{jhagṛo} again
\end{itemize}

\begin{tcolorbox}[breakable,colback=teal!4,colframe=teal!35!black,title={Before you move on}]
What does \mw{झगड़ो} mean? (\textquotedblleft{}A quarrel\textquotedblright{}.) Say it once more.
\end{tcolorbox}
\section[\texorpdfstring{\mw{लड़णो} (laṛṇo)}{लड़णो (laṛṇo)}]{\mw{लड़णो} (laṛṇo) — to fight; to quarrel}

\label{lesson:MW-C244-larno}



\begin{quote}
Before the new one: say the Marwadi for to cheat, then the Marwadi for a quarrel.
\end{quote}

\subsection*{You'll want to know: \mw{लड़णो}}
\textbf{\mw{लड़णो}} — \emph{laṛṇo} — \textquotedblleft{}to fight; to quarrel\textquotedblright{}.

\begin{quote}
\mw{छोरो} \mw{अर} \mw{छोरी} \mw{काल} \mw{लड़्या।}
\end{quote}

\begin{quote}
\emph{chhoro ar chhorī kāl laṛyā.}
\end{quote}

\begin{quote}
The boy and the girl had a fight yesterday.
\end{quote}

Past (he …): \mw{वो} \mw{पाड़ोसी} \mw{लड़्यो}, \emph{vo pāṛosī laṛyo}.

\begin{cousinweb}
\textbf{\mw{लड़णो}} (\emph{laṛṇo}) is inherited from a reconstructed \emph{*laḍ}, \textquotedblleft{}fight\textquotedblright{}. CDIAL gives no Marwari form in this entry; it lists Hindi \emph{laṛnā}, \textquotedblleft{}to fight\textquotedblright{} (CDIAL 10920).
\end{cousinweb}

\begin{grammarlens}[title={What you've built}]
One more word for saying what happened.
\end{grammarlens}

\subsection*{Your turn}
\begin{itemize}
\raggedright
  \item \emph{Say it:} \emph{laṛṇo}
  \item \emph{Say it:} \emph{laṛṇo}, once more
  \item \emph{Say it:} \emph{chhoro ar chhorī kāl laṛyā.} — The boy and the girl had a fight yesterday.
  \item \emph{Recall:} say the Marwadi for to cheat, then the Marwadi for a quarrel, then say \emph{laṛṇo} again
\end{itemize}

\begin{tcolorbox}[breakable,colback=teal!4,colframe=teal!35!black,title={Before you move on}]
What does \mw{लड़णो} mean? (\textquotedblleft{}To fight; to quarrel\textquotedblright{}.) Say it once more.
\end{tcolorbox}
\section[\texorpdfstring{\mw{झूठ} (jhūṭh)}{झूठ (jhūṭh)}]{\mw{झूठ} (jhūṭh) — a lie; a falsehood}

\label{lesson:MW-C244-jhuth}



\begin{quote}
Before the new one: say the Marwadi for a quarrel, then the Marwadi for to fight.
\end{quote}

\subsection*{You'll want to know: \mw{झूठ}}
\textbf{\mw{झूठ}} — \emph{jhūṭh} — \textquotedblleft{}a lie; a falsehood\textquotedblright{}.

\begin{quote}
\mw{झूठ} \mw{मत} \mw{बोलो।}
\end{quote}

\begin{quote}
\emph{jhūṭh mat bolo.}
\end{quote}

\begin{quote}
Don't tell lies.
\end{quote}

A feminine noun.

\begin{cousinweb}
\textbf{\mw{झूठ}} (\emph{jhūṭh}) is inherited from a reconstructed \emph{*jhūṭṭha}, \textquotedblleft{}false\textquotedblright{}; CDIAL lists Marwari \emph{jhūṭh}, \textquotedblleft{}lie\textquotedblright{}, under it (CDIAL 5407).
\end{cousinweb}

\begin{grammarlens}[title={What you've built}]
One more word for saying what happened.
\end{grammarlens}

\subsection*{Your turn}
\begin{itemize}
\raggedright
  \item \emph{Say it:} \emph{jhūṭh}
  \item \emph{Say it:} \emph{jhūṭh}, once more
  \item \emph{Say it:} \emph{jhūṭh mat bolo.} — Don't tell lies.
  \item \emph{Recall:} say the Marwadi for a quarrel, then the Marwadi for to fight, then say \emph{jhūṭh} again
\end{itemize}

\begin{tcolorbox}[breakable,colback=teal!4,colframe=teal!35!black,title={Before you move on}]
What does \mw{झूठ} mean? (\textquotedblleft{}A lie; a falsehood\textquotedblright{}.) Say it once more.
\end{tcolorbox}
\section[\texorpdfstring{\mw{दोस} (dos)}{दोस (dos)}]{\mw{दोस} (dos) — a fault; blame}

\label{lesson:MW-C244-dos}



\begin{quote}
Before the new one: say the Marwadi for to fight, then the Marwadi for a lie.
\end{quote}

\subsection*{You'll want to know: \mw{दोस}}
\textbf{\mw{दोस}} — \emph{dos} — \textquotedblleft{}a fault; blame\textquotedblright{}.

\begin{quote}
\mw{यो} \mw{म्हारो} \mw{दोस} \mw{कोनी।}
\end{quote}

\begin{quote}
\emph{yo mhāro dos konī.}
\end{quote}

\begin{quote}
This is not my fault.
\end{quote}

A masculine noun.

\begin{cousinweb}
\textbf{\mw{दोस}} (\emph{dos}) is inherited from Sanskrit \emph{doṣa}, \textquotedblleft{}fault\textquotedblright{}; CDIAL lists Marwari \emph{dos}, \textquotedblleft{}fault, blame\textquotedblright{}, under it (CDIAL 6587).
\end{cousinweb}

\begin{grammarlens}[title={What you've built}]
One more word for saying what happened.
\end{grammarlens}

\subsection*{Your turn}
\begin{itemize}
\raggedright
  \item \emph{Say it:} \emph{dos}
  \item \emph{Say it:} \emph{dos}, once more
  \item \emph{Say it:} \emph{yo mhāro dos konī.} — This is not my fault.
  \item \emph{Recall:} say the Marwadi for to fight, then the Marwadi for a lie, then say \emph{dos} again
\end{itemize}

\begin{tcolorbox}[breakable,colback=teal!4,colframe=teal!35!black,title={Before you move on}]
What does \mw{दोस} mean? (\textquotedblleft{}A fault; blame\textquotedblright{}.) Say it once more.
\end{tcolorbox}
\section[\texorpdfstring{\mw{धमकी} (dhamkī)}{धमकी (dhamkī)}]{\mw{धमकी} (dhamkī) — a threat}

\label{lesson:MW-C244-dhamki}



\begin{quote}
Before the new one: say the Marwadi for a lie, then the Marwadi for a fault.
\end{quote}

\subsection*{You'll want to know: \mw{धमकी}}
\textbf{\mw{धमकी}} — \emph{dhamkī} — \textquotedblleft{}a threat\textquotedblright{}.

\begin{quote}
\mw{धमकी} \mw{मत} \mw{दो।}
\end{quote}

\begin{quote}
\emph{dhamkī mat do.}
\end{quote}

\begin{quote}
Don't make threats.
\end{quote}

A feminine noun.

\begin{cousinweb}
\textbf{\mw{धमकी}} (\emph{dhamkī}) is inherited from a reconstructed \emph{*dhammakka}, \textquotedblleft{}threat\textquotedblright{}. CDIAL gives no Marwari form in this entry; it lists Hindi and Gujarati \emph{dhamkī} (CDIAL 6736).
\end{cousinweb}

\begin{grammarlens}[title={What you've built}]
One more word for saying what happened.
\end{grammarlens}

\subsection*{Your turn}
\begin{itemize}
\raggedright
  \item \emph{Say it:} \emph{dhamkī}
  \item \emph{Say it:} \emph{dhamkī}, once more
  \item \emph{Say it:} \emph{dhamkī mat do.} — Don't make threats.
  \item \emph{Recall:} say the Marwadi for a lie, then the Marwadi for a fault, then say \emph{dhamkī} again
\end{itemize}

\begin{tcolorbox}[breakable,colback=teal!4,colframe=teal!35!black,title={Before you move on}]
What does \mw{धमकी} mean? (\textquotedblleft{}A threat\textquotedblright{}.) Say it once more.
\end{tcolorbox}
